\documentclass[twocolumn,apj,numberedappendix,appendixfloats]{openjournal}
% Available options:
% [twocolumn] - two-column mode
% [onecolumn] - (default) main text in one-column mode
% [apj]       - typeset in the style of ApJ.
% [apjl]      - (default) typeset in the style of ApJ Letters 
% [tighten]   - some adjustments to approximate grid typesetting
% [numberedappendix]   - number appendix sections as A, B, etc
% [appendixfloats]  - use separate numbering for floats within appendix
% [twocolappendix]  - make appendix in two-col mode in a two-col paper


%%%%% PACKAGES %%%%%
% Optional useful packages
\usepackage{newtxtext,newtxmath,amsmath}
\usepackage[T1]{fontenc}
\usepackage{graphicx}
\usepackage{enumitem}
\usepackage{xcolor}
\usepackage{hyperref}
\usepackage{gensymb}
\usepackage{CJK}
\usepackage{tabularx}
\usepackage{orcidlink}
\hypersetup{
    unicode, 
    colorlinks=true,
    linkcolor=linkcolor,
    citecolor=linkcolor,
    filecolor=linkcolor,
    urlcolor=linkcolor,
}
\definecolor{linkcolor}{rgb}{0.0,0.3,0.5}


% \usepackage{newtxtext,newtxmath,amsmath}
% \usepackage{xcolor}
% \usepackage{textgreek}
% \usepackage[utf8]{inputenc}
% \usepackage[english]{babel}
% \usepackage{hyperref}
% \hypersetup{
%     unicode, 
%     colorlinks=true,
%     linkcolor=linkcolor,
%     citecolor=linkcolor,
%     filecolor=linkcolor,
%     urlcolor=linkcolor,
% }
% \usepackage{color,colortbl}
% \definecolor{linkcolor}{rgb}{0.0,0.3,0.5}
% \usepackage{orcidlink}

%%%%% COMMANDS %%%%%
% Optional useful commands

% I use the following command to indicate comments with \SB{My comment here}.
% Before submission, you should comment this \newcommand out,
% to make sure you do not have any leftover comments!
\newcommand{\SB}[1]{{\textcolor{purple}{SB: #1}}}
\newcommand{\comment}[1]{\textcolor{lightgray}{#1}}

%%%%%%%%%%%%%%%%%%%%%%%%%%%%%%%%%%%%%%%%%%%%%%%%%
%%%%%%%%%%%%%%%%%%%%%%%%%%%%%%%%%%%%%%%%%%%%%%%%%
\begin{document}
\title{The devil is in the detail: Learning from spectrum residuals}

\author{\vspace{-1.5cm}Sven Buder$^{1,*}$\orcidlink{0000-0002-4031-8553}}
\author{David W. Hogg$^{2,3,4}$\orcidlink{0000-0002-4031-8553}}

\affiliation{$^{1}$Research School of Astronomy and Astrophysics, Australian National University, Canberra, ACT 2611, Australia}
\affiliation{$^{2}$Center for Computational Astrophysics, Flatiron Institute, 162 Fifth Avenue, New York, NY 10010, USA}
\affiliation{$^{3}$Max-Planck-Institut für Astronomie, Königstuhl 17, D-69117 Heidelberg, Germany}
\affiliation{$^{4}$Center for Cosmology and Particle Physics, Department of Physics, New York University, 726 Broadway, New York, NY 10003, USA}



\thanks{E-mail: sven.buder@anu.edu.au}
\thanks{$^*$Australian Research Council DECRA Fellow}

\begin{abstract}
Fit a model to the residuals of observed and modelled stellar spectra. Where could this work? "Normal" stars: 1D -> 3D effects. Young stars: Activity (use x-ray luminosity or logRHK), ISM density (via extinction if the velocity is small or we expect the velocity to be small? Or we can shift the spectra to the velocity of the ISM)
The idea would be similar to fitting a binary spectrum to the residuals.
\end{abstract}

% Write your keywords here
\begin{keywords}
    {Choose up to 6 keywords from \href{https://academic.oup.com/DocumentLibrary/mnras/keywords.pdf}{here} and write them like: keyword1 -- keyword2.}
\end{keywords}

\maketitle


\section{Introduction}
\label{sec:introduction}

The paper is structured as follows: Section~\ref{sec:data} describes the data of \dots. Section~\ref{sec:analysis} analyses \dots. We discuss our analysis in the context of our hypotheses in Section~\ref{sec:discussion} and bundles our results into an overarching conclusion in Section~\ref{sec:conclusions}.

%%%%%%%%%%%%%%%%%%%%%%%%%%%%%%%%%%%%%%%%%%%%%%%%%
%%%%%%%%%%%%%%%%%%%%%%%%%%%%%%%%%%%%%%%%%%%%%%%%%
\section{Data} \label{sec:data}

\subsection{Residual spectra} \label{sec:data_residuals}

GALAH DR4 residuals -- one of the data products of GALAH DR4 \citep{Buder2024b}.

Because spectra are reported in their native pixel grid (with a wavelength solution fitted to it), we first interpolated all spectra onto the same wavelength grid.

To do this, for each CCD we identified the highest first pixel wavelength, the median dispersion across all spectra, and the lowest last pixel wavelength\footnote{For CCDs 1-4, this resulted in a common wavelength grid of $4719.0299..0.046..4893.8808$, $5655.5298..0.0547..5863.0182$, $6486.0104..0.0631..6726.4572$, and 
$7689.1261..0.0735..7874.3127$\,\AA.}



Write about how spectra were interpolated onto the same wavelength.

\subsection{Stellar Labels} \label{sec:data_labels}

Selecting stars from GALAH DR4 with \texttt{flag\_sp} < 2048, that is, with sufficient spectral quality (SNR > 10) and estimated stellar parameters. We explicitly do not cut out stars with emission lines or binary signatures.

GALAH DR4 stellar parameters for starters -- to test the 3D implications

Could also be logRHK or X-ray luminosity or extinction?

%%%%%%%%%%%%%%%%%%%%%%%%%%%%%%%%%%%%%%%%%%%%%%%%%
%%%%%%%%%%%%%%%%%%%%%%%%%%%%%%%%%%%%%%%%%%%%%%%%%
\section{Analysis} \label{sec:analysis}

\subsubsection{Test 1: Predicting residuals with Stellar Parameters}


\subsubsection{Test 2: Predicting activity with logRHK or X-ray luminosity}


\subsubsection{Test 3: Predicting ISM density with extinction}


%%%%%%%%%%%%%%%%%%%%%%%%%%%%%%%%%%%%%%%%%%%%%%%%%
%%%%%%%%%%%%%%%%%%%%%%%%%%%%%%%%%%%%%%%%%%%%%%%%%
\section{Discussion} \label{sec:discussion}

%%%%%%%%%%%%%%%%%%%%%%%%%%%%%%%%%%%%%%%%%%%%%%%%%
%%%%%%%%%%%%%%%%%%%%%%%%%%%%%%%%%%%%%%%%%%%%%%%%%
\section{Conclusions} \label{sec:conclusions}


%%%%%%%%%%%%%%%%%%%%%%%%%%%%%%%%%%%%%%%%%%%%%%%%%
%%%%%%%%%%%%%%%%%%%%%%%%%%%%%%%%%%%%%%%%%%%%%%%%%
\section*{Acknowledgments}

We acknowledge the traditional owners of the land on which the ANU stands, the Ngunnawal and Ngambri people. We pay our respects to elders past, and present and are proud to continue their tradition of surveying the night sky and its mysteries to better understand our Universe.

SB acknowledges support from the Australian Research Council under grant numbers DE240100150.

%%%%%%%%%%%%%%%%%%%%%%%%%%%%%%%%%%%%%%%%%%%%%%%%%
%%%%%%%%%%%%%%%%%%%%%%%%%%%%%%%%%%%%%%%%%%%%%%%%%
\section*{Data Availability}

All code to reproduce the analysis and figures can be publicly accessed via \url{https://github.com/svenbuder/residual_cannon}.

%%%%%%%%%%%%%%%%%%%%%%%%%%%%%%%%%%%%%%%%%%%%%%%%%
%%%%%%%%%%%%%%%%%%%%%%%%%%%%%%%%%%%%%%%%%%%%%%%%%
\bibliographystyle{mnras}
\bibliography{bib}

%%%%%%%%%%%%%%%%%%%%%%%%%%%%%%%%%%%%%%%%%%%%%%%%%
%%%%%%%%%%%%%%%%%%%%%%%%%%%%%%%%%%%%%%%%%%%%%%%%%
%\begin{appendix}

%\section{Interesting information that is not vital for the main story of your paper}

%\end{appendix}

\end{document}